\documentclass[11pt]{article}
\usepackage[a4paper,margin=1.05in]{geometry}
\usepackage[T1]{fontenc}
\usepackage{amsmath,amssymb,amsthm,mathtools}
\usepackage{stmaryrd}
\usepackage{mathpartir}
\usepackage{enumitem}
\usepackage{xcolor}
\usepackage[colorlinks=true,linkcolor=blue!55!black,citecolor=blue!55!black,urlcolor=blue!55!black]{hyperref}

%%% Theorem environments
\theoremstyle{plain}
\newtheorem{theorem}{Theorem}[section]
\newtheorem{lemma}[theorem]{Lemma}
\newtheorem{corollary}[theorem]{Corollary}
\newtheorem{proposition}[theorem]{Proposition}
\theoremstyle{definition}
\newtheorem{definition}[theorem]{Definition}
\newtheorem{example}[theorem]{Example}
\newtheorem{assumption}[theorem]{Assumption}
\theoremstyle{remark}
\newtheorem{remark}[theorem]{Remark}

%%% Keywords and syntax
\newcommand{\kw}[1]{\mathsf{#1}}
\newcommand{\kif}{\kw{if}\;}
\newcommand{\kthen}{\;\kw{then}\;}
\newcommand{\kelse}{\;\kw{else}\;}
\newcommand{\nil}{\mathbf{0}}
\newcommand{\ttrue}{\mathsf{tt}}
\newcommand{\ffalse}{\mathsf{ff}}
\newcommand{\seq}{;\;}
\newcommand{\pp}{\mathsf{p}}
\newcommand{\qq}{\mathsf{q}}
\newcommand{\rr}{\mathsf{r}}
\newcommand{\Procs}{\mathcal{P}}
\newcommand{\Var}{\mathit{Var}}
\newcommand{\Val}{\mathit{Val}}
\newcommand{\pn}{\mathrm{pn}}
\newcommand{\eval}[3]{#1 \downarrow_{#2} #3}
\newcommand{\conf}[2]{\langle #1,\, #2\rangle}
\newcommand{\lto}[1]{\xrightarrow{\;#1\;}}
\newcommand{\proj}[2]{\llbracket #1 \rrbracket_{#2}}
\newcommand{\projn}[1]{\llbracket #1 \rrbracket}
\newcommand{\tauat}[1]{\tau\mathord{@}#1}
\newcommand{\cast}[2]{#1 \rightsquigarrow #2}
\newcommand{\upd}[3]{#1[#2 \mapsto #3]}
\newcommand{\rn}[1]{{\normalfont\scshape #1}}
\newcommand{\restr}[2]{#1|_{#2}}
\newcommand{\ifcond}[3]{\kif #1 \kthen #2 \kelse #3}
\newcommand{\sendbr}[4]{#1!#2 \mathbin{\triangleright} (#3,\, #4)}
\newcommand{\recvbr}[3]{#1?(#2,\, #3)}
\newcommand{\pif}[3]{\kif #1 \kthen #2 \kelse #3}
\newcommand{\case}[1]{\medskip\noindent\textbf{Case #1.}\;}
\newcommand{\subcase}[1]{\smallskip\noindent\textit{Subcase #1.}\;}
\newcommand{\Bis}{\mathcal{B}}

\title{A Core Choreographic Language with\\Multicast Knowledge of Choice}
\author{}
\date{\today}

\begin{document}
\maketitle

\begin{abstract}
  We present a small choreographic programming language in the style of Core
  Choreographies, with local computation, point-to-point communication and
  conditionals, but without procedures or recursion.  The language departs from
  Core Choreographies in how it propagates knowledge of choice.  Instead of
  selections and merging, a conditional guarded by process $\pp$ is
  implemented by having $\pp$ \emph{multicast the value of its guard} to every
  other process that participates in the conditional, with every participant,
  including $\pp$, selecting its branch in the same step.  At the network level
  this is realised by two symmetric atomic constructs, \emph{send-and-branch}
  for the guarding process and \emph{receive-and-branch} for the receivers.  We
  give the syntax and semantics of choreographies and of network processes,
  define endpoint projection, and prove that projection is correct in both
  directions.  As a consequence of the design, projection is total (no merging
  can fail), no runtime terms are needed in the choreographic semantics, and
  the correctness theorem is an exact labelled bisimulation with no pruning
  relation.
\end{abstract}

\section*{Overview}

Core Choreographies (CC) implements \emph{knowledge of choice} by
\emph{selections}: after a conditional $\ifcond{\pp.e}{C_1}{C_2}$, the process
$\pp$ must inform every other participant of the branch it took with a label
communication of the form $\pp \to \qq[l]$.  Endpoint projection then uses a
partial \emph{merging} operator to combine the two branches of a process that
has not yet learned the choice, and \emph{projectability} becomes a side
condition that may fail.  Moreover, since the projections of the two branches
have been merged, the correspondence between a choreography and its projection
must be stated up to a \emph{pruning} relation that discards the unreachable
alternatives once a choice is known.

This note works out an alternative.  A conditional is realised by a single
synchronous multicast: $\pp$ evaluates its guard to a boolean $v$, sends $v$
to the set $Q$ of all other processes occurring in either branch, and
\emph{in the same step} every participant, $\pp$ included, continues with the
branch selected by $v$.  Processes that occur in neither branch are not
involved at all.  Because every participant learns the choice before executing
any part of a branch, there is nothing to merge and no dead code to prune.

The network-level constructs that implement this are deliberately symmetric.
The guarding process executes a \emph{send-and-branch}
$\sendbr{Q}{e}{B_1}{B_2}$, which evaluates $e$, multicasts the result to $Q$
and continues as $B_1$ or $B_2$ according to that result; each receiver
executes a \emph{receive-and-branch} $\recvbr{\pp}{B_1}{B_2}$, which waits for
a boolean from $\pp$ and continues accordingly.  Since the sender's own
branching is part of the multicast step rather than a separate local action,
there is no intermediate network state in which the announcement has been
made but the sender has not yet committed.  The choreographic semantics
therefore needs no runtime term for a partially executed conditional, and
the correspondence between choreographies and networks is a step-for-step
labelled bisimulation on \emph{all} configurations.

The remainder is organised as follows: \S\ref{sec:chor} gives the syntax and
semantics of choreographies, \S\ref{sec:net} the syntax and semantics of
network processes, \S\ref{sec:epp} endpoint projection, and
\S\ref{sec:correct} the proof of correctness in both directions.

\section{Choreographies}
\label{sec:chor}

\subsection{Preliminaries}

We fix a finite set $\Procs$ of \emph{process names}, ranged over by
$\pp,\qq,\rr$; sets of process names are ranged over by $Q$.  We also fix a
set $\Var$ of \emph{variables}, ranged over by $x,y,z$, and a set $\Val$ of
\emph{values}, ranged over by $v,w$, containing the booleans $\ttrue$ and
$\ffalse$.

Local \emph{expressions} $e$ are a parameter of the language.  A \emph{local
store} is a total function $s : \Var \to \Val$.  We write
$\eval{e}{s}{v}$ when $e$ evaluates to $v$ in store $s$.

\begin{assumption}\label{ass:eval}
  Evaluation is total and deterministic: for every $e$ and $s$ there is exactly
  one $v$ with $\eval{e}{s}{v}$.  Furthermore, every expression used as the
  guard of a conditional evaluates to a boolean in every store.
\end{assumption}

\noindent
The second half of the assumption is a typing concern that is orthogonal to the
present development; any standard local type system enforces it.  A (global)
\emph{store} is a function $\sigma : \Procs \to (\Var \to \Val)$; we write
$\sigma_\pp$ for $\sigma(\pp)$ and $\upd{\sigma}{\pp.x}{v}$ for the store that
agrees with $\sigma$ except that $\upd{\sigma}{\pp.x}{v}_\pp(x) = v$.

\subsection{Syntax}

\begin{definition}[Choreographies]\label{def:chor-syntax}
  Actions $\eta$ and choreographies $C$ are given by the grammar
  \[
  \begin{array}{r@{\;\;}c@{\;\;}l@{\qquad}l}
    \eta & ::= & \pp.x := e & \text{(local computation)}\\
         & \mid & \pp.e \to \qq.x & \text{(communication, $\pp \neq \qq$)}\\[1ex]
    C & ::= & \nil & \text{(terminated)}\\
      & \mid & \eta \seq C & \text{(action)}\\
      & \mid & \ifcond{\pp.e}{C_1}{C_2} & \text{(conditional)}
  \end{array}
  \]
  For a conditional with branches $C_1, C_2$ we write $C_\ttrue = C_1$ and
  $C_\ffalse = C_2$, so that $C_v$ is the branch selected by the boolean $v$.
\end{definition}

The action $\pp.x := e$ has $\pp$ evaluate $e$ and store the result in its
local variable $x$.  The action $\pp.e \to \qq.x$ has $\pp$ evaluate $e$ and
send the result to $\qq$, which stores it in $x$.  In
$\ifcond{\pp.e}{C_1}{C_2}$ the process $\pp$ evaluates its guard $e$ and the
choreography continues as $C_1$ or $C_2$ accordingly.  As in CC, conditionals
occur only in tail position: there is no sequential composition after a
conditional, and a continuation shared by both branches is written in each of
them.  Unlike CC, there are no selections: knowledge of choice is built into
the semantics of the conditional itself.

\begin{definition}[Process names]\label{def:pn}
  The set $\pn(\cdot)$ of process names occurring in an action or
  choreography is defined by
  \[
  \begin{array}{l@{\quad}l}
    \pn(\pp.x := e) = \{\pp\} &
    \pn(\pp.e \to \qq.x) = \{\pp,\qq\}\\[0.5ex]
    \pn(\nil) = \emptyset &
    \pn(\eta \seq C) = \pn(\eta) \cup \pn(C)\\[0.5ex]
    \multicolumn{2}{l}{
    \pn(\ifcond{\pp.e}{C_1}{C_2}) = \{\pp\}\cup\pn(C_1)\cup\pn(C_2).}
  \end{array}
  \]
  The \emph{receivers} of a conditional at $\pp$ with branches $C_1,C_2$ are
  the processes $(\pn(C_1)\cup\pn(C_2)) \setminus \{\pp\}$.
\end{definition}

Note that the receivers of a conditional are all processes other than $\pp$
that occur in \emph{either} branch, even if they occur in only one of them.
The receiver set may be empty.

\begin{example}\label{ex:chor}
  Let $\Procs = \{\pp,\qq,\rr\}$ and consider
  \[
    C_{\mathrm{ex}} \;=\;
    \pp.x := 2\cdot x \seq
    \ifcond{\pp.(x > 10)}
       {\big(\rr.n := n+1 \seq \pp.x \to \qq.y \seq \qq.y \to \rr.z \seq \nil\big)}
       {\big(\qq.y := 0 \seq \nil\big)}.
  \]
  Process $\pp$ doubles $x$ and then tests whether it exceeds $10$.  In the
  positive case, $\rr$ increments a counter and $\pp$'s value is forwarded to
  $\rr$ via $\qq$; in the negative case, $\qq$ resets $y$.  The receivers of the
  conditional are $\{\qq,\rr\}$; note that $\rr$ occurs only in the first
  branch.
\end{example}

\subsection{Semantics}

The semantics is a labelled transition system over \emph{configurations}
$\conf{C}{\sigma}$.  Transition labels record which processes take part in a
step and what kind of step it is.

\begin{definition}[Labels]\label{def:labels}
  Labels $\mu$ are given by
  \[
    \mu ::= \tauat{\pp} \;\mid\; \pp \to \qq \;\mid\; \cast{\pp}{Q}
  \]
  denoting, respectively, an internal step at $\pp$, a communication from
  $\pp$ to $\qq$, and a multicast from $\pp$ to the set $Q$ (with $\pp\notin Q$).
  The processes involved in a label are
  $\pn(\tauat{\pp}) = \{\pp\}$, $\pn(\pp\to\qq) = \{\pp,\qq\}$ and
  $\pn(\cast{\pp}{Q}) = \{\pp\}\cup Q$.
\end{definition}

\begin{definition}[Choreographic semantics]\label{def:chor-sem}
  The relation $\conf{C}{\sigma} \lto{\mu} \conf{C'}{\sigma'}$ is the least
  relation closed under the rules in Figure~\ref{fig:chor-sem}.
\end{definition}

\begin{figure}[t]
\begin{mathpar}
\inferrule*[lab=\rn{C|Local}]
  {\eval{e}{\sigma_\pp}{v}}
  {\conf{\pp.x := e \seq C}{\sigma} \lto{\tauat{\pp}} \conf{C}{\upd{\sigma}{\pp.x}{v}}}
\and
\inferrule*[lab=\rn{C|Com}]
  {\eval{e}{\sigma_\pp}{v}}
  {\conf{\pp.e \to \qq.x \seq C}{\sigma} \lto{\pp\to\qq} \conf{C}{\upd{\sigma}{\qq.x}{v}}}
\\
\inferrule*[lab=\rn{C|Cast}]
  {\eval{e}{\sigma_\pp}{v} \\ Q = (\pn(C_1)\cup\pn(C_2))\setminus\{\pp\}}
  {\conf{\ifcond{\pp.e}{C_1}{C_2}}{\sigma} \lto{\cast{\pp}{Q}} \conf{C_v}{\sigma}}
\and
\inferrule*[lab=\rn{C|Delay}]
  {\conf{C}{\sigma} \lto{\mu} \conf{C'}{\sigma'} \\ \pn(\eta)\cap\pn(\mu) = \emptyset}
  {\conf{\eta \seq C}{\sigma} \lto{\mu} \conf{\eta \seq C'}{\sigma'}}
\end{mathpar}
\caption{Semantics of choreographies.}
\label{fig:chor-sem}
\end{figure}

Rules \rn{C|Local} and \rn{C|Com} are as in CC.  Rule \rn{C|Cast} executes a
conditional in a single step: the guard is evaluated at $\pp$, the resulting
boolean $v$ is announced to the receivers, and the choreography continues as
the selected branch $C_v$.  The label $\cast{\pp}{Q}$ records the announcement
and its recipients; the store is unchanged.

Rule \rn{C|Delay} implements \emph{out-of-order execution}, which is needed
because a network of independent processes may execute actions of different
processes in any order.  It lets a step of the continuation $C$ happen before
the action $\eta$, provided the processes involved are disjoint.

\begin{remark}[No delay through a conditional]\label{rem:no-delay-cond}
  CC has a rule allowing a step to be delayed past a conditional when both
  branches can perform that same step, which is exactly the situation in
  which merging is needed.  There is no such rule here, and none is needed.
  Every receiver is waiting for $\pp$'s announcement and cannot act before
  it, while any process that is not a receiver does not occur in either branch
  and so has nothing to do.  The only step a configuration
  $\conf{\ifcond{\pp.e}{C_1}{C_2}}{\sigma}$ can take is \rn{C|Cast}.  This is
  confirmed by the network semantics in \S\ref{sec:net} and by the soundness
  proof in \S\ref{sec:correct}.
\end{remark}

\begin{example}\label{ex:run}
  Continuing Example~\ref{ex:chor}, let $\sigma$ be a store with
  $\sigma_\pp(x) = 7$ and $\sigma_\rr(n) = 0$, and abbreviate the two branches
  by $C_1 = \rr.n := n+1 \seq C_1'$ with
  $C_1' = \pp.x \to \qq.y \seq \qq.y \to \rr.z \seq \nil$, and
  $C_2 = \qq.y := 0 \seq \nil$.  One execution of $C_{\mathrm{ex}}$ is
  \begin{align*}
    \conf{C_{\mathrm{ex}}}{\sigma}
      &\lto{\tauat{\pp}} \conf{\ifcond{\pp.(x>10)}{C_1}{C_2}}{\sigma_1}
      && \text{by \rn{C|Local}}\\
      &\lto{\cast{\pp}{\{\qq,\rr\}}} \conf{C_1}{\sigma_1}
      && \text{by \rn{C|Cast}}\\
      &\lto{\tauat{\rr}} \conf{C_1'}{\sigma_2}
      && \text{by \rn{C|Local}}\\
      &\lto{\pp\to\qq} \conf{\qq.y \to \rr.z \seq \nil}{\sigma_3}
      && \text{by \rn{C|Com}}\\
      &\lto{\qq\to\rr} \conf{\nil}{\sigma_4}
      && \text{by \rn{C|Com}}
  \end{align*}
  where $\sigma_1 = \upd{\sigma}{\pp.x}{14}$, $\sigma_2 = \upd{\sigma_1}{\rr.n}{1}$,
  $\sigma_3 = \upd{\sigma_2}{\qq.y}{14}$ and $\sigma_4 = \upd{\sigma_3}{\rr.z}{14}$.
  The third and fourth steps could also be taken in the other order, using
  \rn{C|Delay} to let $\pp$ and $\qq$ communicate before $\rr$ increments its
  counter.
\end{example}

\begin{lemma}[Store locality]\label{lem:chor-locality}
  If $\conf{C}{\sigma} \lto{\mu} \conf{C'}{\sigma'}$ then
  $\sigma'_\rr = \sigma_\rr$ for every $\rr \notin \pn(\mu)$.
\end{lemma}
\begin{proof}
  By induction on the derivation.  \rn{C|Local} changes the store only at
  $\pp \in \pn(\tauat\pp)$, and \rn{C|Com} only at $\qq \in \pn(\pp\to\qq)$.
  \rn{C|Cast} leaves the store unchanged.  \rn{C|Delay} preserves the label
  and the store change of its premise, so the claim follows from the
  induction hypothesis.
\end{proof}

\section{Network Processes}
\label{sec:net}

\subsection{Syntax}

\begin{definition}[Processes and networks]\label{def:net-syntax}
  Behaviours $B$ are given by the grammar
  \[
  \begin{array}{r@{\;\;}c@{\;\;}l@{\qquad}l}
    B & ::= & \nil & \text{(terminated)}\\
      & \mid & x := e \seq B & \text{(local computation)}\\
      & \mid & \qq!e \seq B & \text{(send to $\qq$)}\\
      & \mid & \pp?x \seq B & \text{(receive from $\pp$ into $x$)}\\
      & \mid & \sendbr{Q}{e}{B_1}{B_2} & \text{(send-and-branch: multicast $e$ to $Q$, branch on it)}\\
      & \mid & \recvbr{\pp}{B_1}{B_2} & \text{(receive-and-branch: receive a boolean from $\pp$, branch on it)}
  \end{array}
  \]
  A \emph{network} $N$ is a total function from $\Procs$ to behaviours.  We
  write $\pp[B]$ for the network fragment mapping $\pp$ to $B$, and
  $\pp[B] \mid N$ for the network that maps $\pp$ to $B$ and agrees with $N$
  on all other processes (where $N$ is then understood as a function on
  $\Procs\setminus\{\pp\}$).  Parallel composition $\mid$ is therefore
  commutative and associative.  For a set $S\subseteq\Procs$ we write
  $\restr{N}{S}$ for the restriction of $N$ to $S$, and we say that $N$ and
  $M$ \emph{agree on $S$} if $\restr{N}{S} = \restr{M}{S}$.
\end{definition}

The behaviours $x := e\seq B$, $\qq!e\seq B$ and $\pp?x\seq B$ are the usual
local assignment, send and receive.  The two remaining constructs are the
network-level counterparts of a choreographic conditional and are symmetric
to each other.  The \emph{send-and-branch} $\sendbr{Q}{e}{B_1}{B_2}$ evaluates
$e$ to a boolean, multicasts it to every process in $Q$, and continues as
$B_1$ if the boolean is $\ttrue$ and as $B_2$ if it is $\ffalse$.  The
\emph{receive-and-branch} $\recvbr{\pp}{B_1}{B_2}$ waits for a boolean from
$\pp$ and continues as $B_1$ or $B_2$ accordingly.  There is no separate local
conditional: a send-and-branch with $Q = \emptyset$ evaluates its guard and
branches without communicating, and we may write $\pif{e}{B_1}{B_2}$ as an
abbreviation for $\sendbr{\emptyset}{e}{B_1}{B_2}$.  Both branching constructs
occur only in tail position.  A network is \emph{terminated} if it maps every
process to $\nil$.

\subsection{Semantics}

\begin{definition}[Network semantics]\label{def:net-sem}
  Network configurations are pairs $\conf{N}{\sigma}$ of a network and a
  store.  The relation $\conf{N}{\sigma}\lto{\mu}\conf{N'}{\sigma'}$, with
  labels as in Definition~\ref{def:labels}, is the least relation closed under
  the rules in Figure~\ref{fig:net-sem}.  In rule \rn{N|Com} the processes
  $\pp$ and $\qq$ are distinct; in rule \rn{N|Cast} the processes
  $\pp,\qq_1,\ldots,\qq_n$ are pairwise distinct and $n\geq 0$.
\end{definition}

\begin{figure}[t]
\begin{mathpar}
\inferrule*[lab=\rn{N|Local}]
  {\eval{e}{\sigma_\pp}{v}}
  {\conf{\pp[x := e \seq B] \mid N}{\sigma} \lto{\tauat\pp} \conf{\pp[B]\mid N}{\upd{\sigma}{\pp.x}{v}}}
\and
\inferrule*[lab=\rn{N|Com}]
  {\eval{e}{\sigma_\pp}{v}}
  {\conf{\pp[\qq!e \seq B] \mid \qq[\pp?x \seq B'] \mid N}{\sigma} \lto{\pp\to\qq} \conf{\pp[B]\mid\qq[B']\mid N}{\upd{\sigma}{\qq.x}{v}}}
\\
\inferrule*[lab=\rn{N|Cast}]
  {\eval{e}{\sigma_\pp}{v} \\ Q = \{\qq_1,\ldots,\qq_n\}}
  {\conf{\pp[\sendbr{Q}{e}{B^1}{B^2}] \mid \textstyle\prod_{i=1}^{n}\qq_i[\recvbr{\pp}{B_i^1}{B_i^2}] \mid N}{\sigma}
   \lto{\cast{\pp}{Q}}
   \conf{\pp[B^v] \mid \textstyle\prod_{i=1}^{n}\qq_i[B_i^v] \mid N}{\sigma}}
\end{mathpar}
\caption{Semantics of networks.  In \rn{N|Cast}, $B^v$ denotes $B^1$ if $v=\ttrue$ and $B^2$ if $v=\ffalse$, and likewise for $B_i^v$.}
\label{fig:net-sem}
\end{figure}

Communication is synchronous.  Rule \rn{N|Cast} is a synchronous
one-to-many rendezvous: in a single step the sender evaluates its guard,
every receiver in $Q$ receives the boolean, and the sender and all receivers
select their continuations.  With $Q=\emptyset$ the rule degenerates to a
local conditional at $\pp$, still labelled $\cast{\pp}{\emptyset}$.  No
structural rule for parallel composition is needed because networks are
functions and each rule carries the untouched part of the network as $N$.

\section{Endpoint Projection}
\label{sec:epp}

\begin{definition}[Endpoint projection]\label{def:epp}
  The projection $\proj{C}{\rr}$ of a choreography $C$ onto a process
  $\rr\in\Procs$ is defined by structural recursion on $C$:
  \[
  \begin{array}{l@{\;\;}l@{\quad}l}
    \proj{\nil}{\rr} & = \nil \\[1ex]
    \proj{\pp.x := e \seq C}{\rr} & = x := e \seq \proj{C}{\rr} & \text{if } \rr = \pp\\
                                  & = \proj{C}{\rr} & \text{otherwise}\\[1ex]
    \proj{\pp.e \to \qq.x \seq C}{\rr} & = \qq!e \seq \proj{C}{\rr} & \text{if } \rr = \pp\\
                                       & = \pp?x \seq \proj{C}{\rr} & \text{if } \rr = \qq\\
                                       & = \proj{C}{\rr} & \text{otherwise}\\[1ex]
    \proj{\ifcond{\pp.e}{C_1}{C_2}}{\rr} & = \sendbr{Q}{e}{\proj{C_1}{\pp}}{\proj{C_2}{\pp}} & \text{if } \rr = \pp\\
                                         & = \recvbr{\pp}{\proj{C_1}{\rr}}{\proj{C_2}{\rr}} & \text{if } \rr \in Q\\
                                         & = \nil & \text{otherwise}
  \end{array}
  \]
  where, in the conditional case, $Q = (\pn(C_1)\cup\pn(C_2))\setminus\{\pp\}$
  is the set of receivers.  The projection of $C$ is the network
  $\projn{C}$ with $\projn{C}(\rr) = \proj{C}{\rr}$ for every $\rr\in\Procs$,
  and the projection of a configuration is
  $\projn{C,\sigma} = \conf{\projn{C}}{\sigma}$.
\end{definition}

The cases for $\nil$ and for actions are standard.  A conditional at $\pp$
projects onto $\pp$ as a send-and-branch whose target set is the receivers and
whose continuations are the projections of the two branches onto $\pp$.  It
projects onto a receiver $\rr\in Q$ as a receive-and-branch from $\pp$ whose
two continuations are the projections of the two branches onto $\rr$; if $\rr$
does not occur in one of the branches, the corresponding continuation is
$\nil$.  A process that is neither $\pp$ nor a receiver projects to $\nil$,
because it occurs in neither branch.

\begin{remark}[Totality]\label{rem:total}
  Projection is a total function: every choreography is projectable.  In CC,
  the projection of a conditional onto a process that is not the guard
  requires merging the projections of the two branches, an operation that is
  only defined when the two behaviours agree up to the selections that
  distinguish them; programs in which a process behaves differently in the two
  branches without first being informed by a selection are rejected.  Here the
  announcement of the guard is built into the semantics of the conditional and
  is directed to \emph{every} process that occurs in either branch, so the
  situation that merging is designed to handle never arises.  The price is that
  every participant of a conditional always receives a message, even one whose
  behaviour is identical in both branches.
\end{remark}

\begin{example}\label{ex:epp}
  The projection of $C_{\mathrm{ex}}$ from Example~\ref{ex:chor} is the network
  $\pp[B_\pp] \mid \qq[B_\qq] \mid \rr[B_\rr]$ with
  \[
  \begin{array}{l@{\;=\;}l}
    B_\pp & x := 2\cdot x \seq \sendbr{\{\qq,\rr\}}{(x > 10)}{\qq!x \seq \nil}{\nil}\\
    B_\qq & \recvbr{\pp}{\pp?y \seq \rr!y \seq \nil}{y := 0 \seq \nil}\\
    B_\rr & \recvbr{\pp}{n := n+1 \seq \qq?z \seq \nil}{\nil}.
  \end{array}
  \]
  Process $\rr$ occurs only in the first branch of the conditional, so its
  second continuation is $\nil$.  Reading the execution in
  Example~\ref{ex:run} through the projection, after the multicast step $\pp$
  is at $\qq!x\seq\nil$, $\qq$ is at $\pp?y\seq\rr!y\seq\nil$ and $\rr$ is at
  $n := n+1\seq\qq?z\seq\nil$, which is exactly the projection of $C_1$.
\end{example}

\section{Correctness of Endpoint Projection}
\label{sec:correct}

We prove that a choreography and its projection have exactly the same
labelled transitions, and that the projection of the target of a
choreographic step is the target of the corresponding network step.  No
hypothesis on the configuration is needed.

\subsection{Auxiliary lemmas}

\begin{lemma}[Network locality and framing]\label{lem:frame}
  Suppose $\conf{N}{\sigma}\lto{\mu}\conf{N'}{\sigma'}$.
  \begin{enumerate}[label=(\roman*),nosep]
  \item\label{it:frame-loc} $N'(\rr) = N(\rr)$ and $\sigma'_\rr = \sigma_\rr$ for every $\rr\notin\pn(\mu)$.
  \item\label{it:frame-frame} For every network $M$ that agrees with $N$ on
    $\pn(\mu)$, $\conf{M}{\sigma}\lto{\mu}\conf{M'}{\sigma'}$, where $M'$
    agrees with $N'$ on $\pn(\mu)$ and with $M$ on $\Procs\setminus\pn(\mu)$.
  \end{enumerate}
\end{lemma}
\begin{proof}
  By inspection of Figure~\ref{fig:net-sem}.  In each rule, the processes
  displayed explicitly on the left-hand side are exactly the processes in
  $\pn(\mu)$: $\pp$ for \rn{N|Local}, $\pp$ and $\qq$ for \rn{N|Com}, and
  $\pp,\qq_1,\ldots,\qq_n$ for \rn{N|Cast}.  The rest of the network is passed
  through as $N$ unchanged, and the store is modified only at a process in
  $\pn(\mu)$, which gives~\ref{it:frame-loc}.  For~\ref{it:frame-frame}, the
  premises of each rule mention only the behaviours of the processes in
  $\pn(\mu)$ and the store $\sigma$; if $M$ agrees with $N$ on $\pn(\mu)$ then
  the same rule instance applies to $\conf{M}{\sigma}$, with the remainder
  $\restr{M}{\Procs\setminus\pn(\mu)}$ in place of
  $\restr{N}{\Procs\setminus\pn(\mu)}$, yielding the same label, the same
  behaviours at $\pn(\mu)$ and the same store $\sigma'$.
\end{proof}

\begin{lemma}[Inversion for networks]\label{lem:inversion}
  Suppose $\conf{N}{\sigma}\lto{\mu}\conf{N'}{\sigma'}$ and $\rr\in\pn(\mu)$.
  Then $N(\rr)\neq\nil$, and exactly one of the following holds.
  \begin{enumerate}[label=(\alph*),nosep]
  \item\label{it:inv-local} $N(\rr) = x := e\seq B$; then the rule is \rn{N|Local}, $\mu = \tauat\rr$, $N'(\rr) = B$, and $\sigma' = \upd{\sigma}{\rr.x}{v}$ where $\eval{e}{\sigma_\rr}{v}$.
  \item\label{it:inv-send} $N(\rr) = \qq!e\seq B$; then the rule is \rn{N|Com}, $\mu = \rr\to\qq$, $N(\qq) = \rr?x\seq B'$ for some $x$ and $B'$, $N'(\rr) = B$, $N'(\qq) = B'$, and $\sigma' = \upd{\sigma}{\qq.x}{v}$ where $\eval{e}{\sigma_\rr}{v}$.
  \item\label{it:inv-recv} $N(\rr) = \pp?x\seq B$; then the rule is \rn{N|Com}, $\mu = \pp\to\rr$, $N(\pp) = \rr!e\seq B'$ for some $e$ and $B'$, $N'(\rr) = B$, $N'(\pp) = B'$, and $\sigma' = \upd{\sigma}{\rr.x}{v}$ where $\eval{e}{\sigma_\pp}{v}$.
  \item\label{it:inv-cast} $N(\rr) = \sendbr{Q}{e}{B^1}{B^2}$; then the rule is \rn{N|Cast}, $\mu = \cast{\rr}{Q}$, $N(\qq) = \recvbr{\rr}{B_\qq^1}{B_\qq^2}$ for every $\qq\in Q$, $N'(\rr) = B^v$ and $N'(\qq) = B_\qq^v$ for every $\qq\in Q$, where $\eval{e}{\sigma_\rr}{v}$, and $\sigma' = \sigma$.
  \item\label{it:inv-brecv} $N(\rr) = \recvbr{\pp}{B_1}{B_2}$; then the rule is \rn{N|Cast}, $\mu = \cast{\pp}{Q}$ for some $Q$ with $\rr\in Q$, $N(\pp) = \sendbr{Q}{e}{B^1}{B^2}$ for some $e$, $B^1$, $B^2$, $N'(\rr) = B_v$ where $\eval{e}{\sigma_\pp}{v}$, and $\sigma' = \sigma$.
  \end{enumerate}
\end{lemma}
\begin{proof}
  By inspection of Figure~\ref{fig:net-sem}.  In every rule, each process in
  $\pn(\mu)$ is displayed on the left-hand side with a behaviour of one of the
  five non-terminated shapes, so $N(\rr)\neq\nil$.  The five shapes are
  pairwise distinct, and each shape occurs in exactly one role of exactly one
  rule: assignment only as the subject of \rn{N|Local}, send and receive only
  as the sender and receiver of \rn{N|Com}, and send-and-branch and
  receive-and-branch only as the sender and a receiver of \rn{N|Cast}.  Hence
  the shape of $N(\rr)$ determines the rule and the role of $\rr$ in it, and
  the remaining claims in each item are read off from that rule.
\end{proof}

\begin{lemma}[Projection of prefixes]\label{lem:proj-prefix}
  \begin{enumerate}[label=(\roman*),nosep]
  \item\label{it:pp-act} If $\rr\notin\pn(\eta)$ then $\proj{\eta\seq C}{\rr} = \proj{C}{\rr}$.
  \item\label{it:pp-nil} If $\rr\notin\pn(C)$ then $\proj{C}{\rr} = \nil$.
  \end{enumerate}
\end{lemma}
\begin{proof}
  \ref{it:pp-act} is immediate from Definition~\ref{def:epp}.
  \ref{it:pp-nil} is by induction on $C$: for $\eta\seq C'$ use
  \ref{it:pp-act} and the induction hypothesis; for a conditional at $\pp$
  with receivers $Q$, $\rr\notin\pn(C)$ means $\rr\notin\{\pp\}\cup Q$, so the
  projection is $\nil$ by definition.
\end{proof}

\begin{lemma}[Projection invariance]\label{lem:proj-inv}
  If $\conf{C}{\sigma}\lto{\mu}\conf{C'}{\sigma'}$ and $\rr\notin\pn(\mu)$,
  then $\proj{C'}{\rr} = \proj{C}{\rr}$.
\end{lemma}
\begin{proof}
  By induction on the derivation of the transition.

  \case{\rn{C|Local}}
  $C = \pp.x := e\seq C'$ with $\rr\neq\pp$.  By
  Lemma~\ref{lem:proj-prefix}\ref{it:pp-act}, $\proj{C}{\rr} = \proj{C'}{\rr}$.

  \case{\rn{C|Com}}
  $C = \pp.e\to\qq.x\seq C'$ with $\rr\notin\{\pp,\qq\}$.  As in the previous
  case.

  \case{\rn{C|Cast}}
  $C = \ifcond{\pp.e}{C_1}{C_2}$, $C' = C_v$ and
  $\rr\notin\{\pp\}\cup Q = \pn(C)$.  Then $\proj{C}{\rr} = \nil$ by
  definition, and since $\rr\notin\pn(C_v)$, also $\proj{C_v}{\rr} = \nil$ by
  Lemma~\ref{lem:proj-prefix}\ref{it:pp-nil}.

  \case{\rn{C|Delay}}
  $C = \eta\seq C_0$ and $C' = \eta\seq C_0'$ with
  $\conf{C_0}{\sigma}\lto{\mu}\conf{C_0'}{\sigma'}$.  The induction
  hypothesis gives $\proj{C_0'}{\rr} = \proj{C_0}{\rr}$.  If $\rr\notin\pn(\eta)$
  the claim follows from Lemma~\ref{lem:proj-prefix}\ref{it:pp-act}.  If
  $\rr\in\pn(\eta)$, then by Definition~\ref{def:epp} both $\proj{C}{\rr}$ and
  $\proj{C'}{\rr}$ consist of the same prefix (an assignment, a send or a
  receive, determined by $\eta$ and $\rr$) followed by $\proj{C_0}{\rr}$,
  respectively $\proj{C_0'}{\rr}$, which are equal.
\end{proof}

\subsection{Completeness}

\begin{theorem}[Completeness]\label{thm:completeness}
  If $\conf{C}{\sigma}\lto{\mu}\conf{C'}{\sigma'}$, then
  \[
    \conf{\projn{C}}{\sigma}\lto{\mu}\conf{\projn{C'}}{\sigma'}.
  \]
\end{theorem}
\begin{proof}
  By induction on the derivation of the transition.

  \case{\rn{C|Local}}
  $C = \pp.x := e\seq C'$, $\mu = \tauat\pp$, $\eval{e}{\sigma_\pp}{v}$ and
  $\sigma' = \upd{\sigma}{\pp.x}{v}$.  We have $\proj{C}{\pp} = x := e\seq\proj{C'}{\pp}$
  and $\proj{C}{\rr} = \proj{C'}{\rr}$ for $\rr\neq\pp$.  Rule \rn{N|Local}
  gives $\conf{\projn{C}}{\sigma}\lto{\tauat\pp}\conf{N'}{\sigma'}$ with
  $N'(\pp) = \proj{C'}{\pp}$ and $N'(\rr) = \proj{C}{\rr} = \proj{C'}{\rr}$
  otherwise, so $N' = \projn{C'}$.

  \case{\rn{C|Com}}
  $C = \pp.e\to\qq.x\seq C'$, $\mu = \pp\to\qq$, $\eval{e}{\sigma_\pp}{v}$ and
  $\sigma' = \upd{\sigma}{\qq.x}{v}$.  We have
  $\proj{C}{\pp} = \qq!e\seq\proj{C'}{\pp}$, $\proj{C}{\qq} = \pp?x\seq\proj{C'}{\qq}$
  and $\proj{C}{\rr} = \proj{C'}{\rr}$ otherwise.  Rule \rn{N|Com} gives a
  transition labelled $\pp\to\qq$ to $\conf{\projn{C'}}{\sigma'}$.

  \case{\rn{C|Cast}}
  $C = \ifcond{\pp.e}{C_1}{C_2}$, $C' = C_v$, $\mu = \cast{\pp}{Q}$ with
  $Q = (\pn(C_1)\cup\pn(C_2))\setminus\{\pp\}$, $\eval{e}{\sigma_\pp}{v}$ and
  $\sigma' = \sigma$.  By Definition~\ref{def:epp},
  $\proj{C}{\pp} = \sendbr{Q}{e}{\proj{C_1}{\pp}}{\proj{C_2}{\pp}}$,
  $\proj{C}{\qq} = \recvbr{\pp}{\proj{C_1}{\qq}}{\proj{C_2}{\qq}}$ for $\qq\in Q$,
  and $\proj{C}{\rr} = \nil$ for $\rr\notin\{\pp\}\cup Q$.  Rule \rn{N|Cast}
  applies and gives $\conf{\projn{C}}{\sigma}\lto{\cast{\pp}{Q}}\conf{N'}{\sigma}$
  with $N'(\pp) = \proj{C_v}{\pp}$, $N'(\qq) = \proj{C_v}{\qq}$ for $\qq\in Q$
  and $N'(\rr) = \nil$ for $\rr\notin\{\pp\}\cup Q$.  For the last group,
  $\rr\notin\pn(C_v)$, so $\proj{C_v}{\rr} = \nil = N'(\rr)$ by
  Lemma~\ref{lem:proj-prefix}\ref{it:pp-nil}.  Hence $N' = \projn{C_v} = \projn{C'}$.

  \case{\rn{C|Delay}}
  $C = \eta\seq C_0$ and $C' = \eta\seq C_0'$ with
  $\conf{C_0}{\sigma}\lto{\mu}\conf{C_0'}{\sigma'}$ and
  $\pn(\eta)\cap\pn(\mu) = \emptyset$.  The induction hypothesis gives
  $\conf{\projn{C_0}}{\sigma}\lto{\mu}\conf{\projn{C_0'}}{\sigma'}$.  By
  Lemma~\ref{lem:proj-prefix}\ref{it:pp-act}, $\projn{C}$ agrees with
  $\projn{C_0}$ on $\pn(\mu)$, since $\pn(\mu)$ is disjoint from $\pn(\eta)$.
  By Lemma~\ref{lem:frame}\ref{it:frame-frame},
  $\conf{\projn{C}}{\sigma}\lto{\mu}\conf{N'}{\sigma'}$ where
  $N'(\rr) = \proj{C_0'}{\rr}$ for $\rr\in\pn(\mu)$ and $N'(\rr) = \proj{C}{\rr}$
  otherwise.  It remains to show $N' = \projn{C'}$.  For $\rr\in\pn(\mu)$ we
  have $\rr\notin\pn(\eta)$, so $\proj{C'}{\rr} = \proj{C_0'}{\rr} = N'(\rr)$.
  For $\rr\notin\pn(\mu)$, Lemma~\ref{lem:proj-inv} gives
  $\proj{C_0'}{\rr} = \proj{C_0}{\rr}$, and therefore
  $\proj{\eta\seq C_0'}{\rr} = \proj{\eta\seq C_0}{\rr} = N'(\rr)$, since both
  sides either add the same prefix to equal projections (if $\rr\in\pn(\eta)$)
  or are the equal projections themselves.
\end{proof}

\subsection{Soundness}

\begin{theorem}[Soundness]\label{thm:soundness}
  If $\conf{\projn{C}}{\sigma}\lto{\mu}\conf{N'}{\sigma'}$, then there is a
  choreography $C'$ such that
  \[
    \conf{C}{\sigma}\lto{\mu}\conf{C'}{\sigma'}
    \qquad\text{and}\qquad
    N' = \projn{C'}.
  \]
\end{theorem}
\begin{proof}
  By structural induction on $C$.  Throughout, $N$ denotes $\projn{C}$.

  \case{$C = \nil$}
  Every process of $N$ is $\nil$, so by Lemma~\ref{lem:inversion} no
  transition is possible and there is nothing to prove.

  \case{$C = \eta\seq C_0$}
  We distinguish whether the step involves a process of $\eta$.

  \subcase{$\pn(\eta)\cap\pn(\mu) \neq \emptyset$}
  First let $\eta = \pp.x := e$, so that $\pp\in\pn(\mu)$ and
  $N(\pp) = x := e\seq\proj{C_0}{\pp}$.  By
  Lemma~\ref{lem:inversion}\ref{it:inv-local}, $\mu = \tauat\pp$,
  $\sigma' = \upd{\sigma}{\pp.x}{v}$ with $\eval{e}{\sigma_\pp}{v}$, and
  $N'(\pp) = \proj{C_0}{\pp}$; by Lemma~\ref{lem:frame}\ref{it:frame-loc},
  $N'(\rr) = N(\rr) = \proj{C_0}{\rr}$ for $\rr\neq\pp$
  (Lemma~\ref{lem:proj-prefix}\ref{it:pp-act}).  Hence $N' = \projn{C_0}$,
  and \rn{C|Local} gives $\conf{C}{\sigma}\lto{\tauat\pp}\conf{C_0}{\sigma'}$;
  take $C' = C_0$.

  Now let $\eta = \pp.e\to\qq.x$, so that $N(\pp) = \qq!e\seq\proj{C_0}{\pp}$
  and $N(\qq) = \pp?x\seq\proj{C_0}{\qq}$, and $\pp\in\pn(\mu)$ or
  $\qq\in\pn(\mu)$.  In the first case Lemma~\ref{lem:inversion}\ref{it:inv-send}
  applies to $\pp$; in the second case Lemma~\ref{lem:inversion}\ref{it:inv-recv}
  applies to $\qq$.  Either way we obtain $\mu = \pp\to\qq$,
  $\sigma' = \upd{\sigma}{\qq.x}{v}$ with $\eval{e}{\sigma_\pp}{v}$,
  $N'(\pp) = \proj{C_0}{\pp}$ and $N'(\qq) = \proj{C_0}{\qq}$, and by
  Lemma~\ref{lem:frame}\ref{it:frame-loc} and
  Lemma~\ref{lem:proj-prefix}\ref{it:pp-act}, $N'(\rr) = \proj{C_0}{\rr}$ for
  all other $\rr$.  Hence $N' = \projn{C_0}$, and \rn{C|Com} gives
  $\conf{C}{\sigma}\lto{\pp\to\qq}\conf{C_0}{\sigma'}$; take $C' = C_0$.

  \subcase{$\pn(\eta)\cap\pn(\mu) = \emptyset$}
  By Lemma~\ref{lem:proj-prefix}\ref{it:pp-act}, $\projn{C_0}$ agrees with
  $N$ on $\pn(\mu)$.  By Lemma~\ref{lem:frame}\ref{it:frame-frame},
  $\conf{\projn{C_0}}{\sigma}\lto{\mu}\conf{N''}{\sigma'}$ where $N''$ agrees
  with $N'$ on $\pn(\mu)$ and with $\projn{C_0}$ elsewhere.  The induction
  hypothesis yields $C_0'$ with $\conf{C_0}{\sigma}\lto{\mu}\conf{C_0'}{\sigma'}$
  and $\projn{C_0'} = N''$.  Rule \rn{C|Delay} gives
  $\conf{\eta\seq C_0}{\sigma}\lto{\mu}\conf{\eta\seq C_0'}{\sigma'}$; take
  $C' = \eta\seq C_0'$.  It remains to show $\projn{C'} = N'$.  For
  $\rr\in\pn(\mu)$ we have $\rr\notin\pn(\eta)$, so
  $\proj{C'}{\rr} = \proj{C_0'}{\rr} = N''(\rr) = N'(\rr)$.  For
  $\rr\notin\pn(\mu)$, Lemma~\ref{lem:frame}\ref{it:frame-loc} gives
  $N'(\rr) = N(\rr) = \proj{\eta\seq C_0}{\rr}$, and Lemma~\ref{lem:proj-inv}
  gives $\proj{C_0'}{\rr} = \proj{C_0}{\rr}$, hence
  $\proj{\eta\seq C_0'}{\rr} = \proj{\eta\seq C_0}{\rr} = N'(\rr)$.

  \case{$C = \ifcond{\pp.e}{C_1}{C_2}$}
  Let $Q = (\pn(C_1)\cup\pn(C_2))\setminus\{\pp\}$.  By
  Definition~\ref{def:epp}, $N(\pp) = \sendbr{Q}{e}{\proj{C_1}{\pp}}{\proj{C_2}{\pp}}$,
  $N(\qq) = \recvbr{\pp}{\proj{C_1}{\qq}}{\proj{C_2}{\qq}}$ for $\qq\in Q$, and
  $N(\rr) = \nil$ otherwise.  By Lemma~\ref{lem:inversion} every process in
  $\pn(\mu)$ has a non-terminated behaviour, so $\pn(\mu)\subseteq\{\pp\}\cup Q$.
  We claim $\pp\in\pn(\mu)$.  Otherwise $\pn(\mu)$ contains some $\qq\in Q$,
  whose behaviour is a receive-and-branch from $\pp$, and
  Lemma~\ref{lem:inversion}\ref{it:inv-brecv} gives $\mu = \cast{\pp}{Q'}$ for
  some $Q'$, contradicting $\pp\notin\pn(\mu)$.  So $\pp\in\pn(\mu)$ and, since
  $N(\pp)$ is a send-and-branch, Lemma~\ref{lem:inversion}\ref{it:inv-cast}
  gives $\mu = \cast{\pp}{Q}$, $\sigma' = \sigma$, $N'(\pp) = \proj{C_v}{\pp}$
  and $N'(\qq) = \proj{C_v}{\qq}$ for $\qq\in Q$, where $\eval{e}{\sigma_\pp}{v}$;
  by Lemma~\ref{lem:frame}\ref{it:frame-loc}, $N'(\rr) = \nil$ for
  $\rr\notin\{\pp\}\cup Q$.  Rule \rn{C|Cast} gives
  $\conf{C}{\sigma}\lto{\cast{\pp}{Q}}\conf{C_v}{\sigma}$; take $C' = C_v$.
  Finally $\projn{C_v} = N'$: at $\pp$ and at every $\qq\in Q$ both are
  $\proj{C_v}{\cdot}$ by the above, and at $\rr\notin\{\pp\}\cup Q$ we have
  $\rr\notin\pn(C_v)$, so $\proj{C_v}{\rr} = \nil = N'(\rr)$ by
  Lemma~\ref{lem:proj-prefix}\ref{it:pp-nil}.
\end{proof}

\subsection{Consequences}

Together, Theorems~\ref{thm:completeness} and~\ref{thm:soundness} say that
projection is a functional bisimulation.

\begin{corollary}[Bisimulation]\label{cor:bisim}
  Let $\Bis$ be the relation
  \[
    \Bis = \{\,(\conf{C}{\sigma},\, \conf{\projn{C}}{\sigma}) \mid C \text{ a choreography and } \sigma \text{ a store}\,\}.
  \]
  Then $\Bis$ is a strong labelled bisimulation between the transition
  systems of Definitions~\ref{def:chor-sem} and~\ref{def:net-sem}:
  \begin{enumerate}[label=(\roman*),nosep]
  \item if $(\conf{C}{\sigma}, \conf{N}{\sigma})\in\Bis$ and
    $\conf{C}{\sigma}\lto{\mu}\conf{C'}{\sigma'}$, then
    $\conf{N}{\sigma}\lto{\mu}\conf{N'}{\sigma'}$ for some $N'$ with
    $(\conf{C'}{\sigma'}, \conf{N'}{\sigma'})\in\Bis$;
  \item if $(\conf{C}{\sigma}, \conf{N}{\sigma})\in\Bis$ and
    $\conf{N}{\sigma}\lto{\mu}\conf{N'}{\sigma'}$, then
    $\conf{C}{\sigma}\lto{\mu}\conf{C'}{\sigma'}$ for some $C'$ with
    $(\conf{C'}{\sigma'}, \conf{N'}{\sigma'})\in\Bis$.
  \end{enumerate}
\end{corollary}
\begin{proof}
  (i) is Theorem~\ref{thm:completeness} with $N' = \projn{C'}$, and (ii) is
  Theorem~\ref{thm:soundness}.
\end{proof}

Unlike the corresponding result for CC, the network side of the bisimulation
is the projection itself and not a network related to it by pruning: since
every participant of a conditional learns the outcome in the very step that
resolves it, the projection never contains code for a branch that has
already been ruled out.  Unlike a two-step implementation of the conditional
(Remark~\ref{rem:no-runtime}), the result holds for every configuration,
without any well-formedness hypothesis.

\begin{lemma}[Progress]\label{lem:progress}
  If $C\neq\nil$ then $\conf{C}{\sigma}\lto{\mu}\conf{C'}{\sigma'}$ for some
  $\mu$, $C'$ and $\sigma'$.
\end{lemma}
\begin{proof}
  By cases on $C$, using Assumption~\ref{ass:eval}.  If $C = \pp.x := e\seq C_0$
  or $C = \pp.e\to\qq.x\seq C_0$, then $e$ evaluates to some $v$ in
  $\sigma_\pp$ and \rn{C|Local}, respectively \rn{C|Com}, applies.  If
  $C = \ifcond{\pp.e}{C_1}{C_2}$, then $e$ evaluates to a boolean $v$ in
  $\sigma_\pp$ and \rn{C|Cast} applies.
\end{proof}

\begin{corollary}[Deadlock-freedom by construction]\label{cor:df}
  For every choreography $C$ and store $\sigma$, the configuration
  $\conf{\projn{C}}{\sigma}$ has no transition if and only if $\projn{C}$ is
  terminated, if and only if $C = \nil$.  Consequently, every network
  configuration reachable from the projection of a choreography either has a
  transition or is terminated.
\end{corollary}
\begin{proof}
  If $C = \nil$ then $\projn{C}$ is terminated and, by
  Lemma~\ref{lem:inversion}, has no transition.  If $C\neq\nil$ then
  $\projn{C}$ is not terminated: for $C = \eta\seq C_0$ the projection onto
  the subject of $\eta$ starts with an assignment or a send, and for a
  conditional the projection onto $\pp$ is a send-and-branch.  Moreover, by
  Lemma~\ref{lem:progress} $\conf{C}{\sigma}$ has a transition, and by
  Theorem~\ref{thm:completeness} so does $\conf{\projn{C}}{\sigma}$.  This
  establishes the two equivalences.  For the last claim, by
  Corollary~\ref{cor:bisim} every configuration reachable from
  $\conf{\projn{C}}{\sigma}$ is of the form $\conf{\projn{C'}}{\sigma'}$ for
  some choreography $C'$, to which the equivalences apply.
\end{proof}

\end{document}
